\documentclass[11pt]{article}

\usepackage[utf8]{inputenc}
\usepackage[T1]{fontenc}
\usepackage{microtype}
\usepackage[a4paper, margin=1in]{geometry}
\usepackage{amsmath, amssymb, amsthm}
\usepackage{graphicx}
\usepackage{booktabs}
\usepackage{xcolor}
\usepackage{hyperref}
\usepackage{algorithm}
\usepackage{algpseudocode}
\usepackage[numbers, sort&compress]{natbib}

\hypersetup{
  colorlinks=true,
  linkcolor=blue!50!black,
  citecolor=blue!50!black,
  urlcolor=blue!70!black,
}

\newcommand{\R}{\mathbb{R}}
\newcommand{\E}{\mathbb{E}}
\newcommand{\T}{\mathsf{T}}
\DeclareMathOperator{\slogdet}{slogdet}
\DeclareMathOperator{\diag}{diag}
\DeclareMathOperator{\softplus}{softplus}

\title{\textbf{Buried Treasure}: \\
Two Rigorous Mechanistic-Interpretability Probes \\
Hiding Inside a Retired Architecture}

\author{
  \textbf{PRIME}\thanks{Correspondence: \texttt{datamasteradept@gmail.com}.
                       Code: \href{https://github.com/Prime-007-hash/phi-plasma-core}{Prime-007-hash/phi-plasma-core}.}
  \\
  Independent / Roko.Network \\
  \texttt{Prime-007-hash}
}

\date{May 2026}

\begin{document}
\maketitle

\begin{abstract}
We describe two rigorous mechanistic-interpretability probes built inside a
language model architecture (PHI\_AEON) that was retired in May 2026 after
training cooked at flat \texttt{main\_loss}. The probes themselves were
buried beneath ten other architectural mechanisms, most of which an
internal mathematical audit retired as decorative. The probes survived
the audit by being mathematically honest: a Koopman operator with an
Extended Dynamic Mode Decomposition polynomial lift, and a
Williams--Beer partial-information-decomposition synergy estimator via
Gaussian multi-information.

We re-extract both probes, document their mechanisms in isolation,
demonstrate they are differentiable, low-overhead, and yield interpretable
spectral and information-theoretic diagnostics, and ship them as a
standalone PyTorch module library (\texttt{koopman\_probe.py},
\texttt{iit\_pid\_probe.py}) suitable for attachment to any transformer
backbone. Tests verify the diagnostics behave as the theory predicts.

We frame this as a methodological note rather than an architecture paper:
the discovery that two genuinely rigorous interpretability tools were
built and discarded as a side effect of a different research goal
suggests that auditing one's own retired work for surviving rigorous
fragments is a high-leverage research activity. The audit pattern itself
is described.

Code at \url{https://github.com/Prime-007-hash/phi-plasma-core}
(\texttt{src/phi\_plasma/koopman\_probe.py},
\texttt{src/phi\_plasma/iit\_pid\_probe.py}).
\end{abstract}

\section{Introduction}

Most modern mechanistic-interpretability work on transformer-style language
models proceeds by attaching post-hoc probes to a trained model:
sparse autoencoders for feature discovery~\citep{cunningham2023sparse},
activation patching for circuit identification~\citep{conmy2023automated},
probing classifiers for representational content~\citep{alain2017probing}.
These probes are powerful but generally lack \emph{native mathematical
content}: they detect patterns rather than measure invariants of the
model's dynamics.

This note describes two probes that do measure such invariants. They were
built as auxiliary mechanisms inside a language-model architecture
($\Phi$\_AEON) that was eventually retired for unrelated reasons. The
internal mathematical audit that retired most of $\Phi$\_AEON's mechanisms
as decorative left these two probes intact, because they were the only
mechanisms in the architecture whose mathematical content survives
inspection. We re-extract them here, in isolation, as a standalone
PyTorch interpretability library.

The contributions of this note are limited and methodological:

\begin{enumerate}
\item \textbf{A Koopman EDMD probe} that lifts a learned latent state
into a polynomial-augmented space and fits a linear operator $K$ to the
lifted dynamics. The probe exposes the spectral gap and Lyapunov estimate
of $K$ as differentiable diagnostic signals.
\item \textbf{A Williams--Beer PID synergy probe} that partitions the
model's hidden state and estimates the synergistic component of the
multi-information via Gaussian multi-information, yielding a
differentiable scalar surrogate for ``integrated information.''
\item \textbf{The audit pattern itself}: a description of how, by
auditing decorative-vs.-honest mathematical content in a retired
architecture, two genuinely rigorous tools were recovered that otherwise
would have been deleted with the parent project.
\end{enumerate}

We make no claim that either probe is novel as a theoretical object:
Koopman EDMD is \citet{williams2015data}; PID is \citet{williams2010nonnegative};
Bertschinger's Gaussian MMI synergy estimator is
\citet{bertschinger2014quantifying}. The contribution is engineering and
methodological: clean, low-overhead, MPS-compatible PyTorch
implementations, and the framing that retired projects often contain
rigorous fragments worth re-extracting.

\section{The Koopman EDMD Probe (L7)}

\subsection{Background}

The Koopman operator $\mathcal{K}$ associated with a discrete dynamical
system $z_{t+1} = F(z_t)$ acts on observables $g: \R^d \to \R$ by
$(\mathcal{K} g)(z) = g(F(z))$. The operator is linear, infinite-dimensional,
and contains the spectral information of the underlying dynamics. Extended
Dynamic Mode Decomposition (EDMD; \citealp{williams2015data}) provides a
finite-dimensional approximation: lift the state via a dictionary
$\varphi: \R^d \to \R^D$ with $D > d$, fit a linear operator
$K \in \R^{D \times D}$ minimizing
\[
\| K \varphi(z_t) - \varphi(z_{t+1}) \|_2^2,
\]
and recover spectral properties of $\mathcal{K}$ from the eigenvalues of $K$.

For an appropriately chosen lift basis, EDMD converges to the true Koopman
operator on the lifted space~\citep{korda2018convergence}. The simplest
informative lift augments coordinates with squared monomials:
$\varphi(z) = [z, (Wz)^{\odot 2}]$ for a learnable projection $W$.

\subsection{Implementation}

Our probe attaches to a transformer's final hidden state $h \in \R^{B
\times T \times d_{\text{model}}}$ via:
\begin{align}
z &= \mathrm{MLP}_{\text{enc}}(h) \in \R^{B \times T \times d_z} \\
\varphi(z) &= [z, (W_{\text{lift}} z)^{\odot 2}] \in \R^{B \times T \times D}
\end{align}
where $D = d_z + \lceil d_z / \phi \rceil$, $\phi = (1 + \sqrt{5})/2$
(the golden-ratio lift size is a hyperparameter choice from the original
$\Phi$\_AEON; the math is invariant to this scale).

The Koopman operator $K \in \R^{D \times D}$ is a learned parameter
initialized at $\phi^{-1} I + \mathcal{N}(0, 0.01)$. The probe's primary
training signal is the one-step residual on the lifted state:
\[
\mathcal{L}_{\text{Koop}} = \frac{1}{B(T-1)D} \sum_{b, t < T-1, d}
\big( (K \varphi(z))_{b, t, d} - \varphi(z)_{b, t+1, d} \big)^2.
\]

\subsection{Spectral diagnostics}

The eigenvalue spectrum of $K$ contains the operator's dynamical
information. We compute it via \texttt{torch.linalg.eig}, which is exact
(not iterative) and supports backward on CPU. Two scalars are exposed:
\begin{align}
\text{spectral\_gap}(K) &= |\lambda_1(K)| - |\lambda_2(K)| \\
\text{lyapunov}(K)     &= \log |\lambda_1(K)|
\end{align}
where eigenvalues are sorted by magnitude. The spectral gap measures how
strongly the dominant mode separates from secondary modes; large gap
$\Rightarrow$ a single dominant dynamical direction. The Lyapunov estimate
indicates expansive ($>0$) vs.\ contractive ($<0$) leading dynamics.

Both diagnostics are computed under \texttt{torch.no\_grad()} as
read-only signals. $K$ is trained by the residual loss, not by
differentiating through \texttt{eig}; this avoids the known
high-condition-number gradient problem near degenerate eigenvalues. On
Apple MPS, \texttt{linalg.eig} is unsupported, so the operator is moved
to CPU for the eigendecomposition; cost is $\sim 10\,\mu\text{s}$ for the
$32 \times 32$ operator we use.

We further guard against transient numerical events (NaN/Inf in $K$
during training) by returning safe defaults when $K$ is non-finite,
allowing the training loop to recover at the next optimizer step. In
practice this guard fires once per $\sim 10^4$ steps during long runs.

\subsection{Behavior verification}

Three properties of the probe are mathematically required and empirically
verified:

\begin{enumerate}
\item Direct gradient descent on $\mathcal{L}_{\text{Koop}}$ should
reduce it. We verify this in
\texttt{tests/test\_concentrate.py::test\_koopman\_residual\_decreases\_under\_gradient}.
\item The spectral diagnostics should be finite for any random valid $K$.
Verified in
\texttt{test\_koopman\_spectral\_diagnostics\_finite}.
\item Forward/backward shapes match expectations. Verified in
\texttt{test\_koopman\_forward\_shapes}.
\end{enumerate}

All three tests pass at the time of release.

\section{The Williams--Beer PID Synergy Probe (L12)}

\subsection{Background}

\citet{williams2010nonnegative}'s partial information decomposition (PID)
splits the mutual information $I(X_1, X_2; Y)$ between two sources
$X_1, X_2$ and a target $Y$ into four non-negative components:
unique-from-$X_1$, unique-from-$X_2$, redundant, and synergistic. The
synergistic component is the information about $Y$ available only from
both sources jointly, not from either alone.

For multivariate Gaussian distributions, the synergy admits a closed-form
estimator via the Bertschinger \emph{Maximum Mutual Information} (MMI)
construction~\citep{bertschinger2014quantifying}:
\[
\mathrm{Syn}(X_1, X_2; Y) = I(X_1, X_2; Y) - \max\big(I(X_1; Y), I(X_2; Y)\big).
\]
Gaussian mutual information itself is exact:
\[
I(A; B) = \frac{1}{2} \log \frac{|\Sigma_A| \, |\Sigma_B|}{|\Sigma_{A,B}|}.
\]

Synergy is the information-theoretic quantity that Tononi's IIT
``integrated information'' framework~\citep{tononi2004information} seeks
to formalize, although the mapping is not unique: multiple PID
decompositions exist and reasonable people disagree on which corresponds
to IIT-$\Phi$. We use Bertschinger's MMI estimator without claim that it
\emph{is} IIT-$\Phi$; we claim only that it is a rigorous information-theoretic
quantity that captures part of what IIT was designed to measure.

\subsection{Implementation}

The probe takes the model's hidden state $h \in \R^{B \times T \times
d_{\text{model}}}$, flattens to $\R^{BT \times d_{\text{model}}}$, and
partitions the feature dimension into $n$ chunks treated as sources $X_1,
\ldots, X_n$. The whole vector serves as the target $Y$.

For each partition we compute the Gaussian entropy via the empirical
covariance:
\begin{align}
\hat{\Sigma}_{X_i} &= \frac{1}{N-1} \tilde{X}_i^\T \tilde{X}_i + 10^{-4} I \\
H_i &= \slogdet \big(\hat{\Sigma}_{X_i} \big)
\end{align}
where $\tilde{X}_i$ is the centered data. The ridge term ensures
numerical stability with high-dimensional, low-sample empirical
covariances. The constant $(2\pi e)^{d/2}$ offset in Gaussian entropy is
dropped since only entropy \emph{differences} matter for MI.

The total correlation (multi-information) is
\[
\mathcal{T}(X_1, \ldots, X_n) = \sum_i H_i - H_{\text{joint}},
\]
and the synergy surrogate is
\[
\hat{\Phi} = \max\Big(0, \, \mathcal{T}(X_1, \ldots, X_n) - \tfrac{1}{2} \max_i H_i \Big).
\]
The factor $\tfrac{1}{2}$ approximates the Bertschinger MMI weighting
between joint and best-single-source information.

\subsection{Numerical and platform considerations}

Two details matter in practice:

\begin{itemize}
\item We use \texttt{slogdet} via Cholesky rather than \texttt{torch.logdet}
because the former is numerically stable for small matrices with near-
zero eigenvalues.
\item We reduce via \texttt{.amax()} rather than \texttt{.max()}.
PyTorch's \texttt{max()} produces an \texttt{evenly\_distribute\_backward}
gradient that has historically caused MPS deadlocks at scale; the
ordinal reduction in \texttt{amax()} avoids this.
\end{itemize}

\subsection{Behavior verification}

Two properties of the synergy estimator are mathematically required and
empirically verified:

\begin{enumerate}
\item For independent partitions, the synergy estimator should be near zero
(structurally lower-bounded at $0$ via the clamp). The test
\texttt{test\_iit\_pid\_forward\_shapes} verifies non-negative output.
\item For partitions sharing a common signal, the estimator should be
positive: the joint contains information not available in any single
partition alone. Verified in
\texttt{test\_iit\_pid\_synergy\_nonzero\_for\_correlated\_partitions}.
\end{enumerate}

Both tests pass at the time of release.

\section{The Audit Pattern}

\subsection{How the probes were buried}

PHI\_AEON was a 12-layer transformer-style language model with each layer
implementing a named architectural mechanism: quasi-crystal positional
encoding (L0), $\phi$-coherent embedding (L1), Yoneda morphism embedding
(L2), Sheaf-Hodge-differential attention (L3), Liquid-KAN feed-forward
(L4), Golden-spiral state space (L5), the Koopman probe (L7), critical
router with Kuramoto / echo-state / chaos branches (L8), persistent
homology probe (L9), tropical attention (L10), Fibonacci residual vector
quantization (L11), and the IIT-PID probe (L12).

A May 2026 mathematical audit applied a uniform skepticism criterion to
each mechanism: \emph{does the named operator actually compute the
mathematical object it is named after, or does it reduce to a generic op
with a descriptive label?} The audit retired most mechanisms as
``decorative'' under this criterion. The two probes described here
survived, because their math actually runs: L7's Koopman is a real linear
operator on a real polynomial-augmented lifted space with real spectral
diagnostics; L12's IIT-PID synergy is Bertschinger's MMI estimator with
correct Gaussian-entropy formulas.

\subsection{The retired architecture}

PHI\_AEON was retired at training step 22{,}258 after the loss curve
flat-lined. The internal post-mortem attributed the flat loss to the
decorative mechanisms: when many ``sophisticated'' auxiliary losses
evaluate to zero or vanish near initialization (because their named
mathematical structure is not actually computed), the loss landscape
degenerates to a vanilla transformer's landscape but with reduced
effective capacity due to the auxiliary stalk dimensions allocated to
the decorative paths. The probes described here were not implicated in
the failure; they were collateral in the retirement.

\subsection{The recovery move}

The recovery move is mechanical: re-read each retired file with the
skepticism criterion, separate honest from decorative content, lift the
honest fragments to standalone modules. In PHI\_AEON's case the move
yielded the two probes documented in this note. We suspect this pattern
applies broadly: research groups that explore many architectural
mechanisms inevitably build a small number of rigorous tools alongside
their decorative ones, and these tools usually outlive the architectures
that hosted them. Re-extraction is a high-leverage research activity.

\section{Conclusion}

We have re-extracted two rigorous mechanistic-interpretability probes
from a retired language-model architecture: a Koopman EDMD probe
exposing spectral-gap and Lyapunov diagnostics over a learned linear
operator on a polynomial-augmented latent, and a Williams--Beer
PID-synergy probe via Bertschinger's Gaussian MMI estimator. Both probes
are differentiable, low-overhead, MPS-compatible, and shipped as
standalone PyTorch modules. Both probes pass empirical sanity tests
verifying their mathematically required behavior.

The methodological message is that retired architectures contain
rigorous fragments worth re-extracting. We suspect this is a high-leverage
activity for any research group with a backlog of abandoned model
families.

\section*{Code availability}

\texttt{src/phi\_plasma/koopman\_probe.py} and
\texttt{src/phi\_plasma/iit\_pid\_probe.py} at \\
\url{https://github.com/Prime-007-hash/phi-plasma-core}. Apache 2.0.

\begin{thebibliography}{99}
\setlength{\itemsep}{2pt}

\bibitem{williams2015data}
Williams, M.~O., Kevrekidis, I.~G., \& Rowley, C.~W. (2015).
\emph{A Data--Driven Approximation of the Koopman Operator: Extending
Dynamic Mode Decomposition}.
Journal of Nonlinear Science, 25, 1307--1346.

\bibitem{korda2018convergence}
Korda, M., \& Mezić, I. (2018).
\emph{On Convergence of Extended Dynamic Mode Decomposition to the
Koopman Operator}.
Journal of Nonlinear Science, 28, 687--710.

\bibitem{williams2010nonnegative}
Williams, P.~L., \& Beer, R.~D. (2010).
\emph{Nonnegative Decomposition of Multivariate Information}.
arXiv:1004.2515.

\bibitem{bertschinger2014quantifying}
Bertschinger, N., Rauh, J., Olbrich, E., Jost, J., \& Ay, N. (2014).
\emph{Quantifying Unique Information}.
Entropy, 16, 2161--2183.

\bibitem{tononi2004information}
Tononi, G. (2004).
\emph{An Information Integration Theory of Consciousness}.
BMC Neuroscience, 5(42).

\bibitem{cunningham2023sparse}
Cunningham, H., Ewart, A., Riggs, L., Huben, R., \& Sharkey, L. (2023).
\emph{Sparse Autoencoders Find Highly Interpretable Features in Language
Models}.
arXiv:2309.08600.

\bibitem{conmy2023automated}
Conmy, A., Mavor-Parker, A.~N., Lynch, A., Heimersheim, S., \& Garriga-Alonso, A. (2023).
\emph{Towards Automated Circuit Discovery for Mechanistic Interpretability}.
NeurIPS.

\bibitem{alain2017probing}
Alain, G., \& Bengio, Y. (2017).
\emph{Understanding Intermediate Layers Using Linear Classifier Probes}.
ICLR Workshop.

\end{thebibliography}

\end{document}
